\subsection{Applications of the variational characterization of eigenvalues}

In this section, we consider Hilbert spaces over C.

PROPOSITION 7. --- Let u and v be positive compact endomorphisms of E.

\begin{enumerate}
\def\labelenumi{\alph{enumi})}
\item
  One has \(|\lambda_n(u) - \lambda_n(v)|\leqslant \| u - v\|\) for all \(n\in \mathbf{I}\) ;
\item
  If \(u \leqslant v\) , then \(\lambda_n(u) \leqslant \lambda_n(v)\) for all \(n \in I\) .
\end{enumerate}

Let \(n \in I\) and let F be an n-dimensional vector subspace of E. For all \(x \in F\) , we have

\[
| \langle x | v (x) \rangle - \langle x | u (x) \rangle | \leqslant \| u - v \| \| x \| ^ {2},
\]

therefore the inequalities

\[
\mathrm{R} _ {\mathrm{F}} (v) - \| u - v \| \leqslant \mathrm{R} _ {\mathrm{F}} (u) \leqslant \mathrm{R} _ {\mathrm{F}} (v) + \| u - v \|.
\]

The assertion \(a\) ) follows from this according to proposition 6 of IV, p. 154.

If \(u \leqslant v\) , we have \(\langle x | u(x) \rangle \leqslant \langle x | v(x) \rangle\) for all \(x \in E\) . For every \(n \in I\) and every subspace \(F\) of dimension \(n\) , we therefore have \(R_F(u) \leqslant R_F(v)\) , whence \(\lambda_n(u) \leqslant \lambda_n(v)\) (loc. cit.).

PROPOSITION 8. --- Let u be a positive compact endomorphism of E. Let H be a closed subspace of E and \(i_{\mathrm{H}}: \mathrm{H} \to \mathrm{E}\) the canonical injection. Denote by \(u_{\mathrm{H}}\) the endomorphism \(i_{\mathrm{H}}^{*}ui_{\mathrm{H}}\) of H. It is compact and positive.

\begin{enumerate}
\def\labelenumi{\alph{enumi})}
\item
  One has \(\mathrm{I_H}\subset \mathrm{I_E}\) and \(\lambda_n(u_{\mathrm{H}})\leqslant \lambda_n(u)\) for all \(n\in \mathrm{I_H}\) ;
\item
  If H has finite codimension \(k \in N\) in E, then we have \(I_{H} + k \subset I_{E}\) and \(\lambda_{n+k}(u) \leqslant \lambda_{n}(u_{\mathrm{H}})\) for all \(n \in I_{H}\) .
\end{enumerate}

The endomorphism \(u_{\mathrm{H}}\) is compact (prop. 3 of III, p. 5). It is positive because \(\langle x \mid u_{\mathrm{H}}(x) \rangle = \langle i_{\mathrm{H}}(x) \mid u(i_{\mathrm{H}}(x)) \rangle \geqslant 0\) for all \(x \in \mathrm{H}\) .

Let \(n \in I_{H} \subset I_{E}\) . Let F be a subspace of dimension \(n + 1\) of H adapted to \(u_{H}\) . We therefore have \(\lambda_{n}(u_{\mathrm{H}}) = r_{\mathrm{F}}(u_{\mathrm{H}})\) (prop 5 of IV, p. 153, a)), and since moreover \(r_{\mathrm{F}}(u_{\mathrm{H}}) = r_{\mathrm{F}}(u) \leqslant \lambda_{n}(u)\) (prop. 6 of IV, p. 154), we obtain assertion a).

Suppose that H has codimension \(k \in N\) in E and that \(n \in I_{H}\) . Let F be a subspace of H of dimension n adapted to \(u_{H}\) . Its orthogonal in H is equal to \(H \cap F^{\circ}\) , and it is the orthogonal in E of the subspace \(F + H^{\circ}\) of dimension \(n + k\) . Hence \(n + k \in I_{E}\) and (prop 5 of IV, p. 153, b))

\[
\lambda_{n}(u_{\mathrm{H}}) = \sup_{\substack{x\in \mathrm{H}\cap \mathrm{F}^{\circ}\\ x\neq 0}}\frac{\langle x\mid u_{\mathrm{H}}(x)\rangle}{\|x\|^{2}} = \mathrm{R}_{\mathrm{F} + \mathrm{H}^{\circ}}(u)
\]

whence \(\lambda_{n}(u_{\mathrm{H}})\leqslant \lambda_{n + k}(u)\) (prop. 6 of IV, p. 154).

DEFINITION 3. --- Let F be a Hilbert space and let u be a compact linear map from E into F.

For every integer \(n \in \mathrm{I}_{\mathrm{E}}\) we set \(\alpha_{n}(u) = \lambda_{n}(u^{*} \circ u)\) . Let J be the set of \(n \in \mathrm{I}_{\mathrm{E}}\) such that \(\alpha_{n}(u) > 0\) . The family \((\alpha_{n}(u))_{n \in \mathrm{J}}\) is called the sequence of singular values of \(u\) repeated with multiplicity.

We say that the sequence \((\alpha_{n}(u))_{n\in\mathrm{I}_{\mathrm{E}}}\) is the extended sequence of singular values of u.

The sequence \((\alpha_{n}(u))_{n\in\mathrm{I}_{\mathrm{E}}}\) is well defined since the endomorphism \(u^{*}\circ u\) of E is compact (III, p. 5, prop. 3); it is a decreasing family of positive real numbers since \(u^{*}\circ u\) is positive.

PROPOSITION 9. --- Let F be a Hilbert space and let u be a compact linear map from E into F.

\begin{enumerate}
\def\labelenumi{\alph{enumi})}
\tightlist
\item
  For \(n \in \mathrm{I}_{\mathrm{E}}\) , we have
\end{enumerate}

\[
\alpha_ {n} (u) = \sup _ {\mathrm{F} \in \mathscr {F} _ {n + 1}} \inf _ {x \in \mathrm{F} - \{0 \}} \frac {\| u (x) \|}{\| x \|} = \inf _ {\mathrm{F} \in \mathscr {F} _ {n}} \sup _ {x \in \mathrm{F} ^ {\circ} - \{0 \}} \frac {\| u (x) \|}{\| x \|}.
\]

\begin{enumerate}
\def\labelenumi{\alph{enumi})}
\setcounter{enumi}{1}
\tightlist
\item
  Let J be the set of \(n \in I_{E}\) such that \(\alpha_{n}(u) \neq 0\) . There exist orthonormal families \((e_{n})_{n \in \mathrm{J}}\) in E and \((f_{n})_{n \in \mathrm{J}}\) in F such that for every \(x \in E\) , we have
\end{enumerate}

\[
u (x) = \sum_ {n \in J} \alpha_ {n} (u) \langle e _ {n} | x \rangle f _ {n}.
\]

Since \(\langle x \mid u^{*}u(x)\rangle = \|u(x)\|^{2}\) for all \(x \in E\) , the definitions and prop. 6 of IV, p. 154, a), imply the equality of the first assertion. Let \((e_{n})_{n\in\mathrm{I}_{\mathrm{E}}}\) be an orthonormal family of E satisfying the conclusions of prop. 3 of IV, p. 152 applied to the compact positive endomorphism \(u^{*} \circ u\) of E. Set \(f_{n} = \alpha_{n}^{-1}u(e_{n})\) for \(n \in J\) . Reasoning as in the proof of corollary 2 of IV, p. 150, we obtain assertion c).

COROLLARY. --- Let F be a Hilbert space and u a compact linear map from E into F. Let \(w \in \mathcal{L}(\mathrm{E})\) and \(v \in \mathcal{L}(\mathrm{F})\) . For every \(n \in I_{E}\) , we have \(\alpha_{n}(w \circ u \circ v) \leqslant \|v\|\) \(\|w\|\) \(\alpha_{n}(u)\) .

This is a consequence of assertion \(a\) ) of the preceding proposition.

Remarks. --- 1) The sequence \((\alpha_{n}(u))_{n\in\mathrm{I}_{\mathrm{E}}}\) being decreasing, the set J is either equal to \(\mathrm{I}_{\mathrm{E}}\) , or equal to a segment \(\{0,\ldots,m\}\) in \(\mathrm{I}_{\mathrm{E}}\) where \(m\in\mathbf{N}\) . This latter case holds if and only if \(u\) has finite rank.

\begin{enumerate}
\def\labelenumi{\arabic{enumi})}
\setcounter{enumi}{1}
\item
  Since \(\| u(x)\| = \| |u|(x)\|\) for all \(x\in \mathrm{E}\) (I, p. 139, prop. 10), we have \(\alpha_{n}(u) = \alpha_{n}(|u|)\) for all \(n\in \mathrm{I}_{\mathrm{E}}\) .
\item
  If \(u\) is positive, then \(\alpha_{n}(u)=\lambda_{n}(u)\) for all \(n\in\mathrm{I}_{\mathrm{E}}\) (indeed, we then have \(\alpha_{n}(u)^{2}=\lambda_{n}(u^{*}u)=\lambda_{n}(u^{2})=\lambda_{n}(u)^{2}\) according to prop. 4 of IV, p. 153).
\item
  It is possible that \(\alpha_{n}(v\circ u\circ w)=0\) even if \(\alpha_{n}(u)\) is nonzero; in this case, \(\alpha_{n}(v\circ u\circ w)\) is not a singular value of \(v\circ u\circ w\) .
\end{enumerate}
% semantic-fragments: legacy-input-seam
\relax{}
